% Options for packages loaded elsewhere
\PassOptionsToPackage{unicode}{hyperref}
\PassOptionsToPackage{hyphens}{url}
\PassOptionsToPackage{dvipsnames,svgnames,x11names}{xcolor}
%
\documentclass[
  10pt,
  a4paper,
]{article}
\usepackage{amsmath,amssymb}
\usepackage{iftex}
\ifPDFTeX
  \usepackage[T1]{fontenc}
  \usepackage[utf8]{inputenc}
  \usepackage{textcomp} % provide euro and other symbols
\else % if luatex or xetex
  \usepackage{unicode-math} % this also loads fontspec
  \defaultfontfeatures{Scale=MatchLowercase}
  \defaultfontfeatures[\rmfamily]{Ligatures=TeX,Scale=1}
\fi
\usepackage{lmodern}
\ifPDFTeX\else
  % xetex/luatex font selection
\fi
% Use upquote if available, for straight quotes in verbatim environments
\IfFileExists{upquote.sty}{\usepackage{upquote}}{}
\IfFileExists{microtype.sty}{% use microtype if available
  \usepackage[]{microtype}
  \UseMicrotypeSet[protrusion]{basicmath} % disable protrusion for tt fonts
}{}
\makeatletter
\@ifundefined{KOMAClassName}{% if non-KOMA class
  \IfFileExists{parskip.sty}{%
    \usepackage{parskip}
  }{% else
    \setlength{\parindent}{0pt}
    \setlength{\parskip}{6pt plus 2pt minus 1pt}}
}{% if KOMA class
  \KOMAoptions{parskip=half}}
\makeatother
\usepackage{xcolor}
\usepackage[margin=23mm]{geometry}
\usepackage{color}
\usepackage{fancyvrb}
\newcommand{\VerbBar}{|}
\newcommand{\VERB}{\Verb[commandchars=\\\{\}]}
\DefineVerbatimEnvironment{Highlighting}{Verbatim}{commandchars=\\\{\}}
% Add ',fontsize=\small' for more characters per line
\newenvironment{Shaded}{}{}
\newcommand{\AlertTok}[1]{\textcolor[rgb]{1.00,0.00,0.00}{\textbf{#1}}}
\newcommand{\AnnotationTok}[1]{\textcolor[rgb]{0.38,0.63,0.69}{\textbf{\textit{#1}}}}
\newcommand{\AttributeTok}[1]{\textcolor[rgb]{0.49,0.56,0.16}{#1}}
\newcommand{\BaseNTok}[1]{\textcolor[rgb]{0.25,0.63,0.44}{#1}}
\newcommand{\BuiltInTok}[1]{\textcolor[rgb]{0.00,0.50,0.00}{#1}}
\newcommand{\CharTok}[1]{\textcolor[rgb]{0.25,0.44,0.63}{#1}}
\newcommand{\CommentTok}[1]{\textcolor[rgb]{0.38,0.63,0.69}{\textit{#1}}}
\newcommand{\CommentVarTok}[1]{\textcolor[rgb]{0.38,0.63,0.69}{\textbf{\textit{#1}}}}
\newcommand{\ConstantTok}[1]{\textcolor[rgb]{0.53,0.00,0.00}{#1}}
\newcommand{\ControlFlowTok}[1]{\textcolor[rgb]{0.00,0.44,0.13}{\textbf{#1}}}
\newcommand{\DataTypeTok}[1]{\textcolor[rgb]{0.56,0.13,0.00}{#1}}
\newcommand{\DecValTok}[1]{\textcolor[rgb]{0.25,0.63,0.44}{#1}}
\newcommand{\DocumentationTok}[1]{\textcolor[rgb]{0.73,0.13,0.13}{\textit{#1}}}
\newcommand{\ErrorTok}[1]{\textcolor[rgb]{1.00,0.00,0.00}{\textbf{#1}}}
\newcommand{\ExtensionTok}[1]{#1}
\newcommand{\FloatTok}[1]{\textcolor[rgb]{0.25,0.63,0.44}{#1}}
\newcommand{\FunctionTok}[1]{\textcolor[rgb]{0.02,0.16,0.49}{#1}}
\newcommand{\ImportTok}[1]{\textcolor[rgb]{0.00,0.50,0.00}{\textbf{#1}}}
\newcommand{\InformationTok}[1]{\textcolor[rgb]{0.38,0.63,0.69}{\textbf{\textit{#1}}}}
\newcommand{\KeywordTok}[1]{\textcolor[rgb]{0.00,0.44,0.13}{\textbf{#1}}}
\newcommand{\NormalTok}[1]{#1}
\newcommand{\OperatorTok}[1]{\textcolor[rgb]{0.40,0.40,0.40}{#1}}
\newcommand{\OtherTok}[1]{\textcolor[rgb]{0.00,0.44,0.13}{#1}}
\newcommand{\PreprocessorTok}[1]{\textcolor[rgb]{0.74,0.48,0.00}{#1}}
\newcommand{\RegionMarkerTok}[1]{#1}
\newcommand{\SpecialCharTok}[1]{\textcolor[rgb]{0.25,0.44,0.63}{#1}}
\newcommand{\SpecialStringTok}[1]{\textcolor[rgb]{0.73,0.40,0.53}{#1}}
\newcommand{\StringTok}[1]{\textcolor[rgb]{0.25,0.44,0.63}{#1}}
\newcommand{\VariableTok}[1]{\textcolor[rgb]{0.10,0.09,0.49}{#1}}
\newcommand{\VerbatimStringTok}[1]{\textcolor[rgb]{0.25,0.44,0.63}{#1}}
\newcommand{\WarningTok}[1]{\textcolor[rgb]{0.38,0.63,0.69}{\textbf{\textit{#1}}}}
\setlength{\emergencystretch}{3em} % prevent overfull lines
\providecommand{\tightlist}{%
  \setlength{\itemsep}{0pt}\setlength{\parskip}{0pt}}
\setcounter{secnumdepth}{-\maxdimen} % remove section numbering
\ifLuaTeX
  \usepackage{selnolig}  % disable illegal ligatures
\fi
\usepackage{bookmark}
\IfFileExists{xurl.sty}{\usepackage{xurl}}{} % add URL line breaks if available
\urlstyle{same}
\hypersetup{
  colorlinks=true,
  linkcolor={black},
  filecolor={Maroon},
  citecolor={Blue},
  urlcolor={blue},
  pdfcreator={LaTeX via pandoc}}

\author{}
\date{}

\begin{document}

\section{DDTA BA5-T1 canonical referent naming and controlled-authoring
trial -
R1}\label{ddta-ba5-t1-canonical-referent-naming-and-controlled-authoring-trial---r1}

\textbf{Status:} COMPLETED / PROVISIONAL PASS WITH
CANONICAL-REFERENT-NAMING LOWER-BOUND

\textbf{Repository baseline reviewed:}
\texttt{aa4b785dc72c2fbcd20fd04976e77fca3d07bf25}

\textbf{Executed scope:} BA5-T1 only.

\subsection{1. Trial question}\label{trial-question}

Can DDTA impose one exact canonical name for each named shared project
referent across governed documentation, accepted Base Analysis and
derived shared views, with tool enforcement and no normative synonym
mapping, without confusing lexical naming with BA identity or
restricting method-owned taxonomy?

\subsection{2. Hypothesis under test}\label{hypothesis-under-test}

\begin{Shaded}
\begin{Highlighting}[]
\NormalTok{H5{-}T1}

\NormalTok{For a named shared referent in one governed baseline/naming scope:}
\NormalTok{  exactly one canonical name is normative;}
\NormalTok{  all semantic references use that exact token;}
\NormalTok{  shared derived views preserve it;}
\NormalTok{  aliases/synonyms are not accepted as equivalent input;}
\NormalTok{  tool support enforces rather than interprets equivalence.}
\end{Highlighting}
\end{Shaded}

Example token:

\begin{Shaded}
\begin{Highlighting}[]
\NormalTok{cameraIngresso}
\end{Highlighting}
\end{Shaded}

\subsection{3. Fixed prior contracts}\label{fixed-prior-contracts}

The trial holds BA1-BA4 closed unless a material counterexample forces
the smallest reopen.

Relevant prior invariants:

\begin{itemize}
\tightlist
\item
  BA1: \texttt{BAReferent} identity is independently reusable
  methodology-neutral project meaning; lexical category or convenience
  does not create identity.
\item
  BA2: stable semantic keys are distinct from source wording and display
  labels.
\item
  BA3: wording change does not by itself force referent replacement;
  continuity is
  \texttt{RETAIN\ \textbar{}\ REPLACE\ \textbar{}\ RETIRE}.
\item
  BA4: shared rendering preserves BA meaning; method-owned taxonomy and
  interpretation remain downstream.
\end{itemize}

\subsection{4. Canonical-name propagation
trial}\label{canonical-name-propagation-trial}

Assume governed material introduces one referent under:

\begin{Shaded}
\begin{Highlighting}[]
\NormalTok{cameraIngresso}
\end{Highlighting}
\end{Shaded}

The authoring/integration target is:

\begin{Shaded}
\begin{Highlighting}[]
\NormalTok{Decision / FR / Specialized Requirement / Security Requirement}
\NormalTok{    {-}\textgreater{} cameraIngresso}

\NormalTok{accepted BAReferent rendering}
\NormalTok{    {-}\textgreater{} cameraIngresso}

\NormalTok{human projection}
\NormalTok{    {-}\textgreater{} cameraIngresso}

\NormalTok{flow{-}oriented projection}
\NormalTok{    {-}\textgreater{} [FlowParticipant] cameraIngresso}

\NormalTok{assurance{-}oriented projection}
\NormalTok{    {-}\textgreater{} [AssuranceSubject] cameraIngresso}
\end{Highlighting}
\end{Shaded}

The local projection kinds are intentionally incompatible. The entity
name is not.

Result: \textbf{PASS}.

\subsection{5. Alias negative control}\label{alias-negative-control}

Candidate second-document wording:

\begin{Shaded}
\begin{Highlighting}[]
\NormalTok{entryCamera}
\end{Highlighting}
\end{Shaded}

when it intends the same shared referent as \texttt{cameraIngresso}.

If accepted as a second normative identifier, every consumer must either
maintain synonym equivalence or guess identity. That reintroduces
lexical resolution before the semantic architecture is fully integrated.

Disposition: \textbf{REJECTED}.

The valid options are:

\begin{enumerate}
\def\labelenumi{\arabic{enumi}.}
\tightlist
\item
  use \texttt{cameraIngresso}; or
\item
  perform a governed rename for a later baseline.
\end{enumerate}

\subsection{6. Exact-token control}\label{exact-token-control}

The following are not silently equivalent:

\begin{Shaded}
\begin{Highlighting}[]
\NormalTok{cameraIngresso}
\NormalTok{CameraIngresso}
\NormalTok{camera\_ingresso}
\NormalTok{camera{-}ingresso}
\end{Highlighting}
\end{Shaded}

The registry stores the accepted token and validation is exact under the
project naming convention.

Disposition: \textbf{PASS / exact-token lower bound retained}.

\subsection{7. Collision control}\label{collision-control}

Suppose two different referents are authored with:

\begin{Shaded}
\begin{Highlighting}[]
\NormalTok{cameraIngresso}
\end{Highlighting}
\end{Shaded}

within the same naming scope.

The registry can no longer resolve source/projection references
unambiguously.

Disposition: \textbf{REJECTED}.

Therefore canonical names are unique within the active baseline/naming
scope.

\subsection{8. Human projection control}\label{human-projection-control}

A human view wants a friendlier label:

\begin{Shaded}
\begin{Highlighting}[]
\NormalTok{Entry Camera}
\end{Highlighting}
\end{Shaded}

If this text replaces the referent name, the no-synonym invariant is
broken.

Disposition: \textbf{REJECTED AS REFERENT RENAME}.

The view may instead write explanatory prose around the canonical token:

\begin{Shaded}
\begin{Highlighting}[]
\NormalTok{Entrance capture device: cameraIngresso}
\end{Highlighting}
\end{Shaded}

where only \texttt{cameraIngresso} is the project entity identifier.

\subsection{9. Method projection
control}\label{method-projection-control}

A flow-oriented method wants:

\begin{Shaded}
\begin{Highlighting}[]
\NormalTok{FlowParticipant}
\end{Highlighting}
\end{Shaded}

and an assurance-oriented method wants:

\begin{Shaded}
\begin{Highlighting}[]
\NormalTok{AssuranceSubject}
\end{Highlighting}
\end{Shaded}

Both are allowed as projection-owned kinds:

\begin{Shaded}
\begin{Highlighting}[]
\NormalTok{FlowParticipant(cameraIngresso)}
\NormalTok{AssuranceSubject(cameraIngresso)}
\end{Highlighting}
\end{Shaded}

The method taxonomy does not rename the shared referent.

Disposition: \textbf{PASS}.

\subsection{10. Method-owned aggregation
control}\label{method-owned-aggregation-control}

A method creates an item:

\begin{Shaded}
\begin{Highlighting}[]
\NormalTok{external{-}ingress{-}exposure}
\end{Highlighting}
\end{Shaded}

from several traced BA elements including \texttt{cameraIngresso}.

This label is allowed because it names a method-owned
interpretation/aggregation, not the project referent. Its trace still
names/reaches \texttt{cameraIngresso} canonically.

Disposition: \textbf{PASS WITH BA4 BOUNDARY}.

\subsection{11. Governed rename falsification
control}\label{governed-rename-falsification-control}

The strongest test intentionally changes the canonical name:

\begin{Shaded}
\begin{Highlighting}[]
\NormalTok{B0: cameraIngresso}
\NormalTok{B1: cameraNord}
\end{Highlighting}
\end{Shaded}

Assume the independently reusable project meaning is materially
unchanged; only the governed name changes.

If BA5 equated name and identity, the rename would force a new
\texttt{BAReferent}, contradicting BA3's closed identity rule that
wording changes do not by themselves force replacement.

The correct result is:

\begin{Shaded}
\begin{Highlighting}[]
\NormalTok{B0 BAReferent R}
\NormalTok{  canonicalName cameraIngresso}

\NormalTok{BACrossBaselineContinuity}
\NormalTok{  R@B0 {-}\textgreater{} RETAIN {-}\textgreater{} R@B1}

\NormalTok{B1 BAReferent R}
\NormalTok{  canonicalName cameraNord}
\end{Highlighting}
\end{Shaded}

Historical B0 projections remain \texttt{cameraIngresso}. Rebuilt B1
shared projections use \texttt{cameraNord}.

Disposition: \textbf{PASS WITH REFINEMENT}.

This falsifies a globally immutable-name hypothesis and closes the
narrower rule:

\begin{quote}
canonical name is exact and invariant within one governed
baseline/naming scope, while governed cross-baseline rename may retain
semantic identity.
\end{quote}

\subsection{12. Source authority
control}\label{source-authority-control}

Negative pattern:

\begin{Shaded}
\begin{Highlighting}[]
\NormalTok{ThreatForge sees entryCamera}
\NormalTok{  {-}\textgreater{} silently decides it means cameraIngresso}
\NormalTok{  {-}\textgreater{} writes BA/project binding}
\end{Highlighting}
\end{Shaded}

Disposition: \textbf{REJECTED}.

Allowed pattern:

\begin{Shaded}
\begin{Highlighting}[]
\NormalTok{tool detects unregistered entryCamera}
\NormalTok{  {-}\textgreater{} blocks or proposes cameraIngresso}
\NormalTok{  {-}\textgreater{} author/governance accepts correction}
\NormalTok{  {-}\textgreater{} governed material uses cameraIngresso}
\end{Highlighting}
\end{Shaded}

The tool enforces the registry; it does not become the registry's
semantic authority.

\subsection{13. BA reopen pressure}\label{ba-reopen-pressure}

The trial tries to force additional semantic machinery.

\subsubsection{New BAE family?}\label{new-bae-family}

Not needed. Naming metadata attaches to/resolves for
\texttt{BAReferent}.

\subsubsection{\texorpdfstring{New BA2 operator such as
\texttt{nameOf}?}{New BA2 operator such as nameOf?}}\label{new-ba2-operator-such-as-nameof}

Not needed. Canonical naming is lexical/governance metadata, not a
methodology-neutral project assertion requiring proposition identity.

\subsubsection{BA3 reopen?}\label{ba3-reopen}

Not needed. Existing continuity already handles a retained identity
under wording/name change.

\subsubsection{BA4 reopen?}\label{ba4-reopen}

Not needed. BA5 simply narrows shared referent rendering: local
projection kinds remain free, shared referent names become
exact-canonical.

\subsubsection{Documentation metamodel
reopen?}\label{documentation-metamodel-reopen}

Not forced. The constraint is cross-document controlled authoring and
naming governance rather than a new structural field family for
Decision/FR/SR/Security Requirement.

\subsection{14. Trial dispositions}\label{trial-dispositions}

\begin{Shaded}
\begin{Highlighting}[]
\NormalTok{canonical name exact within baseline/scope                 PASS}
\NormalTok{same shared referent same name across governed docs        PASS}
\NormalTok{same shared referent same name across shared views         PASS}
\NormalTok{alias/synonym as second normative identifier               REJECTED}
\NormalTok{case/format variant as implicit alias                      REJECTED}
\NormalTok{collision between two referents                            REJECTED}
\NormalTok{canonicalName as semantic identity                         REJECTED}
\NormalTok{governed cross{-}baseline rename with RETAIN                 PASS}
\NormalTok{projection{-}owned type vocabulary                           PASS}
\NormalTok{method{-}owned interpretation label                          PASS}
\NormalTok{tool registry validation/completion                        PASS}
\NormalTok{tool autonomous alias normalization                        REJECTED}
\NormalTok{new BAE family                                             NOT FORCED}
\NormalTok{new BA2 operator                                           NOT FORCED}
\NormalTok{BA1{-}BA4 reopen                                             NOT TRIGGERED}
\NormalTok{BA5                                                        STARTED / NOT CLOSED}
\end{Highlighting}
\end{Shaded}

\subsection{15. Result}\label{result}

BA5-T1 passes with one important refinement.

The initial strong controlled-vocabulary direction survives, but
\texttt{canonicalName} cannot be treated as an eternal identity key. The
accepted provisional lower bound is:

\begin{Shaded}
\begin{Highlighting}[]
\NormalTok{same named referent}
\NormalTok{+ same governed baseline/naming scope}
\NormalTok{    {-}\textgreater{} one exact canonicalName everywhere shared}

\NormalTok{cross{-}baseline governed rename}
\NormalTok{    {-}\textgreater{} may RETAIN BAReferent identity}
\NormalTok{    {-}\textgreater{} new baseline uses new canonicalName consistently}
\end{Highlighting}
\end{Shaded}

No synonym registry is required by the current evidence.

\subsection{16. Next smallest falsification
target}\label{next-smallest-falsification-target}

BA5-T2 should pressure-test \textbf{registry coverage and governed
extension}, not general synonym normalization:

\begin{itemize}
\tightlist
\item
  operators;
\item
  operator-scoped roles;
\item
  semantic kinds such as \texttt{store} and \texttt{contract};
\item
  controlled values;
\item
  new canonical-term admission; and
\item
  tool enforcement/completion across the complete registry.
\end{itemize}

Only concrete failure should force alias/synonym machinery.

\end{document}
